\section*{Hoofdstuk \ref{ch:b1:layered-earth} --- De gelaagde aarde}

\begin{solution}{exo:b1:layered-earth:1}
De volumefracties zijn de kubussen van de stralenverhoudingen: kern
$(3480/6371)^3 = 0.1630$, dat is \qty{16.3}{\percent}; \omterm{def:b1:layered-earth:layers}{binnenkern}
$(1221.5/6371)^3 = 0.00705$, \qty{0.70}{\percent}.
\end{solution}

\begin{solution}{exo:b1:layered-earth:2}
Boven: $\Phi = 13.7166^2 - \tfrac43\times7.26466^2 = 188.145 - 70.367 =
\qty{117.78}{km^2.s^{-2}}$, $K = \rho\Phi = \qty{5566.45}{kg.m^{-3}}\times
\qty{117.78e6}{m^2.s^{-2}} = \qty{655.6}{GPa}$. Onder: $\Phi = 8.06482^2 =
\qty{65.04}{km^2.s^{-2}}$, $K = 9903.49\times\num{65.04e6} = \qty{644.1}{GPa}$.
De compressiemodulus verandert nauwelijks (\qty{-1.7}{\percent}). Wat verandert is de
schuifmodulus, $\mu = \rho v_S^2 = \qty{293.8}{GPa}$ boven en $0$ onder (een vloeistof), en de
dichtheid, die met \qty{78}{\percent} stijgt: vandaar de daling van $v_P$ ondanks een bijna gelijke $K$.
\end{solution}

\begin{solution}{exo:b1:layered-earth:3}
$\sin i_c = 6.2/8.1 = 0.7654$, $i_c = \ang{49.95}$. Kritieke afstand
$2\times30\times\tan\ang{49.95} = 60\times1.1906 = \qty{71.4}{km}$. \omterm{def:b1:layered-earth:head-wave}{Kruisafstand}
$60\sqrt{14.3/1.9} = 60\times2.743 = \qty{164.6}{km}$.
\end{solution}

\begin{solution}{exo:b1:layered-earth:4}
$g = Gm/r^2 = \num{6.674e-11}\times\num{1.94e24}/(\num{3.48e6})^2 =
\qty{10.69}{m.s^{-2}}$, \qty{9}{\percent} meer dan aan het oppervlak. Naar de kern toe daalt het kwadraat van de straal
tot $(3480/6371)^2 = 0.298$ van zijn oppervlaktewaarde, terwijl de aantrekkende massa slechts daalt tot
$1.94/5.972 = 0.325$ van het totaal, omdat de kern een derde van de massa bevat: $0.325/0.298 = 1.09$.
\end{solution}

\begin{solution}{exo:b1:layered-earth:5}
Nabij de bron $v_1 = 60/10.0 = \qty{6.0}{km.s^{-1}}$. Veraf
$v_2 = (300 - 200)/(45.8 - 33.3) = 100/12.5 = \qty{8.0}{km.s^{-1}}$; de
intercepttijd is $t_i = 33.3 - 200/8.0 = \qty{8.3}{s}$. Dan $\sin i_c =
0.75$, $\cos i_c = 0.6614$, en $h = t_i v_1/(2\cos i_c) = 8.3\times6.0/1.3229
= \qty{37.6}{km}$. Controle: $x_c = 2\times37.6\times\sqrt{14/2} = \qty{199}{km}$,
waar beide golven aankomen op $199/6.0 = 199/8.0 + 8.3 = \qty{33.2}{s}$.
\end{solution}

\begin{solution}{exo:b1:layered-earth:6}
Voor uniforme $\rho$ is $I = \tfrac{8\pi}{3}\rho\int_0^R r^4\,\mathrm dr =
\tfrac{8\pi}{15}\rho R^5$ en $M = \tfrac43\pi\rho R^3$, dus
$I/MR^2 = (8/15)/(4/3) = 0.4$. Twee lagen ($r_c = \qty{3.48e6}{m}$,
$R = \qty{6.371e6}{m}$): $M = \tfrac43\pi\bigl(11000\,r_c^3 + 4450\,(R^3 -
r_c^3)\bigr) = \qty{5.977e24}{kg}$, \qty{0.08}{\percent} boven de gemeten
massa; $I = \tfrac{8\pi}{15}\bigl(11000\,r_c^5 + 4450\,(R^5 - r_c^5)\bigr)$ geeft
$I/MR^2 = 0.3457$. De massa klopt, maar de factor is \qty{4.5}{\percent}
te hoog: de werkelijke \omterm{def:b1:layered-earth:layers}{mantel} is niet uniform, haar dichtheid neemt met de diepte toe, wat massa naar binnen verplaatst.
\end{solution}

\begin{solution}{exo:b1:layered-earth:7}
$\Phi = 12.16^2 - \tfrac43\times6.673^2 = 147.87 - 59.37 =
\qty{88.49}{km^2.s^{-2}}$. Met de diepte $z$ is
$\mathrm d\rho/\mathrm dz = \rho g/\Phi = 4860\times9.9/\num{88.49e6} =
\qty{5.44e-4}{kg.m^{-4}}$, dat is \qty{0.544}{kg.m^{-3}} per kilometer, of
\qty{0.54}{g.cm^{-3}} per \qty{1000}{km}. PREM stijgt over deze \qty{1000}{km} met $5.357 - 4.844 =
0.513$, \qty{6}{\percent} minder: de gradiënt die bovenaan wordt berekend, wordt op het hele interval toegepast, terwijl $\Phi$ met de diepte toeneemt en de gradiënt afneemt.
\end{solution}

\begin{solution}{exo:b1:layered-earth:8}
$P \approx \bar\rho g z = 3570\times9.9\times\num{670e3} = \qty{23.7}{GPa}$
(PREM: \qty{23.9}{GPa}). Het is de druk, ongeveer \qty{24}{GPa}, waarbij ringwoodiet uiteenvalt in
bridgmaniet en ferroperiklaas: de discontinuïteit op \qty{660}{km}.
\end{solution}

\begin{solution}{exo:b1:layered-earth:9}
(a)~Vloeibaar ijzer: $10.90\times1.05 = 11.4$ tot $10.90\times1.10 =
\qty{12.0}{g.cm^{-3}}$. (b)~$12.764/12.166 = 1.049$: een sprong van
\qty{4.9}{\percent} in PREM, iets meer dan de \qty{4.5}{\percent} van recente
studies. (c)~Stollen alleen zou de vaste stof maar weinig dichter maken dan de vloeistof; de grotere sprong betekent dat de
vaste stof minder lichte elementen bevat dan de vloeistof: die geven de voorkeur aan de vloeistof, en het stollen op de \omterm{def:b1:layered-earth:discontinuity}{binnenkerngrens} laat het meeste ervan achter.
\end{solution}

\begin{solution}{exo:b1:layered-earth:10}
Laat $x_* = 2h\tan i_c$. Dan is $\sqrt{x_*^2 + 4h^2} = 2h/\cos i_c$ en
$t_r(x_*) = 2h/(v_1\cos i_c)$. De \omterm{def:b1:layered-earth:head-wave}{kopgolf} geeft
$t(x_*) = 2h\tan i_c/v_2 + 2h\cos i_c/v_1 = \frac{2h}{v_1}\bigl(\sin^2 i_c/\cos
i_c + \cos i_c\bigr) = 2h/(v_1\cos i_c)$, met $\sin i_c = v_1/v_2$: hetzelfde.
De helling van de hyperbool is $t_r'(x) = x/(v_1\sqrt{x^2 + 4h^2})$; bij $x_*$ is zij
$\sin i_c/v_1 = 1/v_2$, de helling van de lijn van de \omterm{def:b1:layered-earth:head-wave}{kopgolf}. Omdat
$t_r'' > 0$, is de hyperbool convex en ligt zij boven elk van haar raaklijnen: voor
$x > x_*$ geldt $t_r > t_{\text{kop}}$. Op de kritieke afstand raakt de reflectie de grens precies onder $i_c$ en ontstaat de \omterm{def:b1:layered-earth:head-wave}{kopgolf};
daarna komt de \omterm{def:b1:layered-earth:head-wave}{kopgolf}, die wint doordat zij in het snelle medium loopt, het eerst aan.
\end{solution}

\begin{solution}{exo:b1:layered-earth:11}
$m(r) = \tfrac43\pi\rho r^3$, dus $g = Gm/r^2 = \tfrac43\pi G\rho r$. Dan is
$P(r) = \int_r^R \rho\,\tfrac43\pi G\rho s\,\mathrm ds = \tfrac23\pi G\rho^2(R^2
- r^2)$. In het middelpunt is $P = \tfrac23\pi\times\num{6.674e-11}\times5513^2\times
(\num{6.371e6})^2 = \qty{172}{GPa}$, \qty{47}{\percent} van de \qty{363.85}{GPa} van PREM: het concentreren van de massa naar het midden verdubbelt de centrale druk ongeveer.
\end{solution}

\begin{solution}{exo:b1:layered-earth:12}
Met $z$ de diepte is $\mathrm d\rho/\mathrm dz = \rho g/\Phi$; is $g/\Phi$ een constante $a$, dan is $\mathrm d(\ln\rho)/\mathrm dz = a$ en $\rho = \rho_0
e^{a(z - z_0)}$. Hier is $a(z - z_0) = 7.9\times\num{2258.5e3}/\num{87e6} =
0.2051$ en $\rho = 9.90\times e^{0.2051} = 9.90\times1.2276 =
\qty{12.15}{g.cm^{-3}}$, \qty{0.2}{\percent} onder PREM's $12.17$: de hele dichtheidstoename door de \omterm{def:b1:layered-earth:layers}{buitenkern} is zelfcompressie.
\end{solution}

\begin{solution}{pb:b1:layered-earth:1}
\textbf{1.} $V = \tfrac43\pi(\num{6.371e6})^3 = \qty{1.0832e21}{m^3}$.

\textbf{2.} $\tfrac43\pi(r_2^3 - r_1^3)$: \omterm{def:b1:layered-earth:layers}{korst}, bovenste \omterm{def:b1:layered-earth:layers}{mantel} en overgangszone
\qty{3.0706e20}{m^3}; onderste \omterm{def:b1:layered-earth:layers}{mantel} \qty{5.9961e20}{m^3}; \omterm{def:b1:layered-earth:layers}{buitenkern}
\qty{1.6890e20}{m^3}; \omterm{def:b1:layered-earth:layers}{binnenkern} \qty{7.634e18}{m^3}.

\textbf{3.} \qty{28.35}{\percent}, \qty{55.36}{\percent}, \qty{15.59}{\percent}
en \qty{0.70}{\percent}: de onderste \omterm{def:b1:layered-earth:layers}{mantel} is de grootste schil, meer dan de helft
van het volume van de aarde.

\textbf{4.} $15.59 + 0.70 = \qty{16.3}{\percent}$.

\textbf{5.} \omterm{def:b1:layered-earth:layers}{Korst}, bovenste \omterm{def:b1:layered-earth:layers}{mantel} en overgangszone $\num{3.0706e20}\times3570 =
\qty{1.0962e24}{kg}$; onderste \omterm{def:b1:layered-earth:layers}{mantel} \qty{2.9381e24}{kg}; \omterm{def:b1:layered-earth:layers}{buitenkern}
\qty{1.8410e24}{kg}; \omterm{def:b1:layered-earth:layers}{binnenkern} \qty{9.841e22}{kg}.

\textbf{6.} \qty{5.9737e24}{kg}, \qty{0.03}{\percent} boven de gemeten
\qty{5.972e24}{kg}.

\textbf{7.} $\num{5.9737e24}/\num{1.0832e21} = \qty{5515}{kg.m^{-3}}$.

\textbf{8.} $\num{9.841e22}/\num{5.9737e24} = \qty{1.65}{\percent}$.

\textbf{9.} $(1.0962 + 2.9381)/5.9737 = \qty{67.5}{\percent}$.

\textbf{10.} De kern is \qty{16.3}{\percent} van het volume en
\qty{32.5}{\percent} van de massa, twee keer zoveel: haar gemiddelde dichtheid,
\qty{10986}{kg.m^{-3}}, is ongeveer $2.5$ keer die van de \omterm{def:b1:layered-earth:layers}{mantel}.

\textbf{11.} Een schil van uniforme dichtheid geeft
$I = \tfrac{8\pi}{3}\rho\int_{r_1}^{r_2} r^4\,\mathrm dr = \tfrac{8\pi}{15}\rho
(r_2^5 - r_1^5)$.

\textbf{12.} \omterm{def:b1:layered-earth:layers}{Korst}, bovenste \omterm{def:b1:layered-earth:layers}{mantel} en overgangszone \qty{2.676e37}{kg.m^2}; onderste \omterm{def:b1:layered-earth:layers}{mantel}
\qty{4.525e37}{kg.m^2}; \omterm{def:b1:layered-earth:layers}{buitenkern} \qty{9.27e36}{kg.m^2}; \omterm{def:b1:layered-earth:layers}{binnenkern}
\qty{5.87e34}{kg.m^2}.

\textbf{13.} De onderste \omterm{def:b1:layered-earth:layers}{mantel}: zij heeft de helft van de massa, op grote stralen. De \omterm{def:b1:layered-earth:layers}{binnenkern}
bevat \qty{1.6}{\percent} van de massa, en die ligt allemaal dicht bij de as: haar aandeel is \qty{0.07}{\percent}.

\textbf{14.} $I = \qty{8.1345e37}{kg.m^2}$ en $I/MR^2 = \num{8.1345e37}/
(\num{5.9737e24}\times(\num{6.371e6})^2) = 0.3355$.

\textbf{15.} \qty{1.4}{\percent} boven $0.3308$, en ver onder $0.4$: het model heeft de juiste concentratie van de massa naar het midden, bijna.

\textbf{16.} Ja. Binnen elke schil is de werkelijke dichtheid onderaan hoger dan bovenaan; een uniforme schil van dezelfde massa legt te veel
van die massa dicht bij haar buitenstraal, waar zij met $r^2$ meetelt, en overschat dus $I$.

\textbf{17.} $g = GM/R^2 = \num{6.674e-11}\times\num{5.9737e24}/
(\num{6.371e6})^2 = \qty{9.822}{m.s^{-2}}$.

\textbf{18.} De massa van de kern is $\num{1.8410e24} + \num{9.841e22} =
\qty{1.9394e24}{kg}$, en $g = G m/r_c^2 = \qty{10.69}{m.s^{-2}}$.

\textbf{19.} $(\num{1.0962e24} + \num{2.9381e24})/(\num{3.0706e20} +
\num{5.9961e20}) = \num{4.0343e24}/\num{9.0667e20} = \qty{4450}{kg.m^{-3}}$.

\textbf{20.} $P \approx 4450\times10\times\num{2.891e6} = \qty{128.6}{GPa}$,
\qty{5}{\percent} onder de \qty{135.75}{GPa} van PREM.

\textbf{21.} $\rho_c = \num{1.9394e24}/\num{1.7653e20} = \qty{10986}{kg.m^{-3}}$,
en $\Delta P = \tfrac23\pi\times\num{6.674e-11}\times10986^2\times
(\num{3.48e6})^2 = \qty{204.3}{GPa}$.

\textbf{22.} $128.6 + 204.3 = \qty{332.9}{GPa}$, \qty{8.5}{\percent} onder de
\qty{363.85}{GPa} van PREM. Beide schattingen negeren de toename van de dichtheid met de diepte binnen elk gebied, die meer massa dieper plaatst,
de zwaartekracht naar beneden hoger houdt en $\rho g$, de drukgradiënt, verhoogt.

\textbf{23.} De totale massa (vraag 6), de \omterm{def:b1:comparative-planetology:moi}{traagheidsmomentfactor} (vraag 15) en de zwaartekracht aan het oppervlak (vraag 17, tegen de gemeten
gemiddelde \qty{9.82}{m.s^{-2}}) worden gecontroleerd aan metingen van de aarde; de zwaartekracht in de kern en de drukken kunnen alleen met PREM worden vergeleken,
dat zelf een model is.

\textbf{24.} De kern bevat \textbf{\qty[text-series-to-math=true]{32.5}{\percent}} van de massa van de aarde.
\end{solution}
